% ============================================================
% 附录：Lean 4 辅助形式化验证
% ============================================================
\section{Lean 4 辅助形式化验证}
\label{app:lean4}

本文使用 Lean 4 辅助验证三参数九点修正族六阶 local consistency 不可达性反证中的代数核心。
验证对象不是完整 PDE 收敛定理，而是 local Fourier symbol 展开后得到的有限个多项式系数条件。

\subsection{验证目标}

考虑定理~\ref{thm:6th_impossible}中的三参数九点修正族：
\begin{equation}
  (-L_h - \alpha h^2 L_h^{(5)})u + \sigma(R_h+\beta h^2 I)u = (R_h+\gamma h^2 I)f,
\end{equation}
其中 $\alpha,\beta,\gamma$ 是与 Fourier mode 无关的统一常数。
在 local Fourier symbol 展开中，六阶一致性要求 $h^2$ 项系数 $c_2$ 和
$h^4$ 项系数 $c_4$ 对所有连续波数同时为零。Lean 4 辅助验证的目标是检查这些
系数条件中的有限代数核心：
\begin{enumerate}[itemsep=1pt]
  \item \(c_2=0\) 推出的参数必要关系；
  \item 代入这些必要关系后，两个 Fourier 测试模式对应的 \(c_4=0\) 条件互相矛盾；
  \item 因此不存在统一实数参数使这些必要条件同时成立。
\end{enumerate}

\subsection{验证定理与策略}

表~\ref{tab:lean4_theorems} 列出了 Lean 4 中验证的核心实数域定理。
该表只对应反证链条中的有限个代数命题，不表示从差分模板到 PDE 全局误差估计的完整形式化。
\begin{table}[H]
\centering
\caption{Lean 4 辅助验证的代数核心定理清单}
\label{tab:lean4_theorems}
\resizebox{\textwidth}{!}{%
\begin{tabular}{llll}
\toprule
\textbf{Lean 定理名} & \textbf{数学含义} & \textbf{证明策略} & \textbf{状态} \\
\midrule
\texttt{c4\_num\_mode\_10\_real} & \(c_4(1,0,\gamma)=60\gamma+3\) 的代数展开 & \texttt{ring} & 通过 \\
\texttt{c4\_num\_mode\_11\_real} & \(c_4(1,1,\gamma)=120\gamma-4\) 的代数展开 & \texttt{ring} & 通过 \\
\texttt{sixth\_order\_impossible\_real} & 两个测试模式的 \(c_4=0\) 条件不相容 & \texttt{ring}+\texttt{linarith} & 通过 \\
\texttt{c2\_poisson\_forces\_neg\_gamma} & Poisson 情形下 \(c_2=0\) 推出必要关系 \(\alpha=-\gamma\) & \texttt{linarith} & 通过 \\
\texttt{poisson\_no\_sixth\_order\_real} & Poisson 三参数族的 local symbol 六阶条件不相容 & \texttt{ring}+\texttt{linarith} & 通过 \\
\texttt{c2\_helmholtz\_forces\_neg\_gamma} & Helmholtz 情形下 \(c_2=0\) 推出必要关系 \(\alpha=-\gamma\) & \texttt{linarith} & 通过 \\
\texttt{helmholtz\_no\_sixth\_order\_real} & Helmholtz 三参数族的 local symbol 六阶条件不相容 & \texttt{ring}+\texttt{linarith} & 通过 \\
\bottomrule
\end{tabular}%
}
\end{table}

表中定理只对应反证中的代数步骤。它们说明：一旦 \(c_2=0\) 给出必要参数关系，
两个测试波数 \((1,0)\) 与 \((1,1)\) 的 \(c_4=0\) 条件就无法由同一个常数
\(\gamma\) 同时满足。
这里的 \((1,0)\) 与 \((1,1)\) 是 local Fourier symbol 中的连续波数测试点，
不是 Dirichlet DST-I 离散模态编号。
其中，\texttt{c2\_helmholtz\_forces\_neg\_gamma} 直接验证的是
Helmholtz 情形下 \(c_2=0\) 所需参数关系中的 \(\alpha=-\gamma\) 代数核心；
若在论文推导中使用 \(\beta=\gamma\)，该关系来自相应系数方程的纸面推导，
而不是该 Lean 定理的直接结论。

\subsection{验证范围与局限性}

\textbf{已验证的内容}：
\begin{enumerate}[itemsep=1pt]
  \item \(c_4\) 系数在测试波数 \((1,0)\) 和 \((1,1)\) 上的显式代数表达式；
  \item 两个测试波数的 \(c_4=0\) 条件互相矛盾；
  \item \(c_2=0\) 推出相应参数必要关系的代数步骤；
  \item Poisson 与 Helmholtz 两种情形中 \(c_2 \to c_4\) 反证路径的代数不相容性。
\end{enumerate}

\textbf{未覆盖的内容}：
\begin{enumerate}[itemsep=1pt]
  \item Fourier symbol 从差分模板导出的全过程；
  \item Taylor 展开建模为 \(c_2,c_4\) 的全过程；
  \item Dirichlet DST-I 谱理论或 Neumann/Mixed 边界下的 DCT 理论；
  \item 边界闭合分析；
  \item 完整 PDE 全局误差估计；
  \item Python 求解器、实验脚本或数值实验结果的程序正确性。
\end{enumerate}

\subsection{形式化验证工具链}

\begin{itemize}[itemsep=2pt]
  \item \textbf{Lean 版本}：leanprover/lean4:v4.29.1
  \item \textbf{Mathlib 版本}：v4.29.1（预编译 olean）
  \item \textbf{构建命令}：\texttt{lake build}
  \item \textbf{验证状态}：自有形式化文件无 \texttt{sorry}、\texttt{admit}、\texttt{axiom} 或 \texttt{unsafe}，全部定理通过编译。
\end{itemize}

\begin{remark}
  Lean 4 验证覆盖的是 Fourier symbol 展开之后的有限多项式条件之间的不相容性，
  为论文中的代数反证提供机器可检查的支持。
  它不声称证明完整 PDE 收敛理论、边界闭合理论、DST/DCT 谱理论或 Python 程序正确性。
\end{remark}
